\documentclass[11pt]{article}

% Answers and hints for the practice problems of Lecture 5
\usepackage[margin=1in]{geometry}
\usepackage[T1]{fontenc}
\usepackage{lmodern}
\usepackage{amsmath,amssymb,amsthm}
\usepackage{enumitem}
\usepackage{xr-hyper}
\usepackage{hyperref}
% Problem numbers are read from practice5.aux, so compile practice5.tex first.
\externaldocument{practice5}[practice5.pdf]
\usepackage{xcolor}

\title{Lecture 5: Practice Problems --- Answers and Hints}
\author{Math for Machine Learning}
\date{\today}

\newcommand{\norm}[1]{\lVert #1 \rVert}
\newcommand{\E}{\mathbb{E}}
\DeclareMathOperator{\Var}{Var}
\DeclareMathOperator{\Cov}{Cov}
\DeclareMathOperator{\Corr}{Corr}
\DeclareMathOperator{\KL}{KL}
\DeclareMathOperator{\diag}{diag}
\DeclareMathOperator{\tr}{tr}
\DeclareMathOperator{\Div}{div}
\DeclareMathOperator{\softmax}{softmax}

\begin{document}

\maketitle

\section*{Ordinary Differential Equations}

\begin{description}[leftmargin=0pt, itemsep=0.3em]
    \item[\ref{pr:scalar-ode}.] (a) $x(t) = 3e^{-2t}$, $x(1) = 3e^{-2} \approx 0.406$; it halves after $\frac{\log 2}{2} \approx 0.347$. (b) With $u = 4e^{-t}$: $x' = \frac{u}{(1 + u)^2}$ and $x(1 - x) = \frac{1}{1 + u}\cdot\frac{u}{1 + u}$, which agree, and $x(0) = \frac{1}{5}$. So $x(t) = \sigma(t - \log 4)$, $x(2) \approx 0.649$, and $x(t) \to 1$.
    \item[\ref{pr:existence}.] (a) $x' = \frac{1}{(1 - t)^2} = x^2$ and $x(0) = 1$. The solution tends to $\infty$ as $t \to 1$: it blows up in finite time. The theorem does not apply, since $F(x) = x^2$ is not Lipschitz continuous on $\mathbb{R}$: $|x^2 - y^2| = |x + y||x - y|$, and $|x + y|$ is unbounded. (b) For $x = t^2$: $x' = 2t = 2\sqrt{t^2}$. The function $2\sqrt{|x|}$ is not Lipschitz continuous near $0$, where its slope is unbounded, so uniqueness can fail.
\end{description}

\section*{Linear Systems}

\begin{description}[leftmargin=0pt, itemsep=0.3em]
    \item[\ref{pr:linear-system}.] (a) Eigenvalue $1$ with eigenvector $(1, 1)$, and $-3$ with $(1, -1)$; $(1, 0) = \frac{1}{2}(1, 1) + \frac{1}{2}(1, -1)$, so $\mathbf{x}(t) = \frac{1}{2}e^t(1, 1)^\top + \frac{1}{2}e^{-3t}(1, -1)^\top$ and $\mathbf{x}(1) \approx (1.384, 1.334)^\top$. (b) Exactly the multiples of $(1, -1)$, i.e., the line $x_2 = -x_1$; all other solutions grow along $(1, 1)$. (c) $e^{At} = \frac{1}{2}\begin{pmatrix} e^t + e^{-3t} & e^t - e^{-3t} \\ e^t - e^{-3t} & e^t + e^{-3t} \end{pmatrix}$.
    \item[\ref{pr:rotation-ode}.] (a) $\frac{d}{dt}\norm{\mathbf{x}}^2 = 2\mathbf{x}^\top A\mathbf{x} = 0$, and $e^{At}$ is the rotation by the angle $\omega t$, so the period is $\frac{2\pi}{\omega}$. (b) $\norm{\mathbf{x} + hA\mathbf{x}}^2 = \norm{\mathbf{x}}^2 + 2h\mathbf{x}^\top A\mathbf{x} + h^2\norm{A\mathbf{x}}^2 = (1 + h^2\omega^2)\norm{\mathbf{x}}^2$. One period takes $\frac{2\pi}{0.1} \approx 63$ steps, so the norm grows by $(1.01)^{31.4} \approx 1.37$.
\end{description}

\section*{Gradient Flow and Euler's Method}

\begin{description}[leftmargin=0pt, itemsep=0.3em]
    \item[\ref{pr:gradient-flow}.] (a) $x_1' = -2x_1$ and $x_2' = -8x_2$, so $\mathbf{x}(t) = (2e^{-2t}, e^{-8t})$ and $f(\mathbf{x}(t)) = 4e^{-4t} + 4e^{-16t}$, which decreases. (b) $\mathbf{x}_{k+1} = \mathbf{x}_k - h\nabla f(\mathbf{x}_k) = \big((1 - 2h)x_{1,k}, (1 - 8h)x_{2,k}\big)$, which is gradient descent with $\eta = h$. It converges if $|1 - 2h| < 1$ and $|1 - 8h| < 1$, i.e., $0 < h < 0.25$. For $h = 0.2$, the factors are $0.6$ and $-0.6$ (convergence with oscillation in $x_2$); for $h = 0.3$, $x_2$ is multiplied by $-1.4$ and diverges, although the exact flow converges.
    \item[\ref{pr:euler-error}.] $0.75^4 \approx 0.316$ and $0.875^8 \approx 0.344$, compared with $e^{-1} \approx 0.368$; the errors $0.051$ and $0.024$ differ by a factor of about $2$, as expected for a method with error proportional to $h$.
\end{description}

\section*{Flows, Volumes and Densities}

\begin{description}[leftmargin=0pt, itemsep=0.3em]
    \item[\ref{pr:volume-flow}.] (a) $\Div F = \tr A = -2$, so areas are multiplied by $e^{-2t}$. The eigenvalues are $1$ and $-3$ (triangular matrix), and $\det e^{At} = e^te^{-3t} = e^{-2t}$. (b) $DF = \diag(1 - 3x^2, -1)$, so $\Div F = -3x^2$. $F = \mathbf{0}$ at $(0, 0)$ (divergence $0$) and $(\pm 1, 0)$ (divergence $-3$). Areas are preserved to first order on the line $x = 0$ and compressed everywhere else.
    \item[\ref{pr:continuity-1d}.] (a) $P(x(t) \leq a) = P(x(0) \leq e^ta)$; differentiating in $a$ gives $e^t\varphi(e^ta)$. (b) With $\varphi'(u) = -u\varphi(u)$: $\frac{\partial p_t}{\partial t} = e^t\varphi(e^tx) + e^{2t}x\,\varphi'(e^tx)$, and $-\frac{\partial}{\partial x}(-xp_t) = p_t + x\frac{\partial p_t}{\partial x} = e^t\varphi(e^tx) + xe^{2t}\varphi'(e^tx)$. (c) $\log p_t(x(t)) = t + \log\varphi(x(0))$, which increases at the rate $1$. The flow contracts lengths by $e^{-t}$ (here $\Div F = -1$), so the same probability is squeezed into a smaller interval, and the density increases.
\end{description}

\section*{Stochastic Differential Equations}

\begin{description}[leftmargin=0pt, itemsep=0.3em]
    \item[\ref{pr:brownian}.] (a) $\mathcal{N}(0, 3)$; $B(4) \sim \mathcal{N}(0, 4)$, so $P(B(4) > 2) = P(Z > 1) \approx 0.159$. (b) $\Cov = \min(1, 2) = 1$, so $\Corr = \frac{1}{\sqrt{2}} \approx 0.707$. $B(1) + B(2) = 2B(1) + (B(2) - B(1))$ is a sum of independent Gaussians, $\mathcal{N}(0, 4 + 1) = \mathcal{N}(0, 5)$. (c) $B(2) = 1 + (B(2) - B(1)) \sim \mathcal{N}(1, 1)$, since the increment is independent of $B(1)$. The vector $(B(1), B(2))$ has the covariance matrix $\left(\begin{smallmatrix} 1 & 1 \\ 1 & 2 \end{smallmatrix}\right)$, so $B(1) \mid B(2) = 2$ is $\mathcal{N}(\frac{1}{2} \cdot 2, \, 1 - \frac{1}{2}) = \mathcal{N}(1, \frac{1}{2})$: the mean lies on the straight line from $(0, 0)$ to $(2, 2)$ (a ``Brownian bridge'').
    \item[\ref{pr:euler-maruyama}.] (a) $X_1 = 1 - 0.2 + 0.5\sqrt{0.1} \cdot 0.5 \approx 0.879$ and $X_2 = 0.8 \cdot 0.879 - 0.5\sqrt{0.1} \approx 0.545$. (b) $\E[X_{k+1}] = 0.8\,\E[X_k]$, so $\E[X_k] = 0.8^k$, and $\E[X_{10}] \approx 0.107$, compared with $e^{-2} \approx 0.135$: the error of Euler's method for $x' = -2x$. (c) The SDE: $\frac{\sigma^2}{2\theta} = \frac{0.25}{4} = 0.0625$. The chain: $\frac{\sigma^2}{\theta(2 - \theta h)} = \frac{0.25}{3.6} \approx 0.0694$. A stationary distribution exists if $|1 - 2h| < 1$, i.e., $0 < h < 1$.
    \item[\ref{pr:sqrt-h}.] The variance is $\sigma^2h^{2\gamma} \cdot \frac{t}{h} = \sigma^2t\,h^{2\gamma - 1}$. For $\gamma > \frac{1}{2}$, the noise disappears in the limit (we get an ODE); for $\gamma < \frac{1}{2}$, it explodes. Only $\gamma = \frac{1}{2}$ gives a limit with variance $\sigma^2t$, independently of $h$, which is why Brownian motion has $\Var(B(t)) = t$ and the Euler--Maruyama noise is scaled by $\sqrt{h}$.
    \item[\ref{pr:quadratic-variation}.] (a) The increments are $\sqrt{h}\,Z_k$ with i.i.d.\ $Z_k \sim \mathcal{N}(0, 1)$, and $hZ_k^2$ has the mean $h$ and the variance $h^2(\E[Z^4] - 1) = 2h^2$. Summing $n$ independent terms gives $\E[Q_n] = nh = T$ and $\Var(Q_n) = 2nh^2 = 2Th \to 0$, so $Q_n \to T$ by Chebyshev's inequality. (b) $|f((k + 1)h) - f(kh)| = |f'(\xi_k)|h \leq Ch$ with $C = \max|f'|$, so the sum is at most $nC^2h^2 = C^2Th \to 0$. Since $Q_n \to T > 0$, Brownian paths cannot be continuously differentiable.
\end{description}

\section*{Langevin Dynamics and Diffusion Models}

\begin{description}[leftmargin=0pt, itemsep=0.3em]
    \item[\ref{pr:ula}.] (a) $x_{k+1} = x_k - \frac{h}{s^2}(x_k - \mu) + \sqrt{2h}\,z_k$, so $a = 1 - \frac{h}{s^2}$. (b) If $x_k - \mu \sim \mathcal{N}(0, v)$, then $x_{k+1} - \mu \sim \mathcal{N}(0, a^2v + 2h)$, and $v = a^2v + 2h$ gives $v = \frac{2h}{1 - a^2} = \frac{2h}{2h/s^2 - h^2/s^4} = \frac{2s^2}{2 - h/s^2}$; the chain converges to it if $|a| < 1$, i.e., $0 < h < 2s^2$. The values give $\frac{0.5}{1.96} \approx 0.2551$ instead of $0.25$; the bias vanishes as $h \to 0$. (c) Gradient descent on $U$, which converges to the minimizer $\mu$. For a posterior, the minimizer of $U = -\log p(\theta \mid \mathcal{D})$ is the MAP estimate; the noise turns optimization into sampling.
    \item[\ref{pr:diffusion-forward}.] (a) $x_t$ is Gaussian as a sum of independent Gaussians. The mean is multiplied by $\sqrt{1 - \beta}$ in each step, and the variance satisfies $v_t = (1 - \beta)v_{t-1} + \beta$; if $v_{t-1} = 1 - (1 - \beta)^{t-1}$, then $v_t = (1 - \beta) - (1 - \beta)^t + \beta = 1 - (1 - \beta)^t$. So $x_t = (1 - \beta)^{t/2}x_0 + \sqrt{1 - (1 - \beta)^t}\,z$ can be sampled in one step, without simulating the chain. (b) $t = 10$: mean $2.71$, variance $0.183$; $t = 100$: $1.09$ and $0.867$; $t = 500$: $0.019$ and $1.000$. For large $t$, $x_t$ is approximately $\mathcal{N}(0, 1)$ for every data point; generation starts from this distribution and runs the process backwards. (c) Euler--Maruyama gives $x_{k+1} = x_k - \frac{h}{2}x_k + \sqrt{h}\,z_k = (1 - \frac{\beta}{2})x_k + \sqrt{\beta}\,z_k$ for $h = \beta$. Its stationary variance $\frac{\sigma^2}{2\theta} = 1$ agrees with (b).
\end{description}

\end{document}
